\chapter*{Sobre el autor}

Carlos de Anta Puig es economista con más de veinticinco años de experiencia profesional en microeconometría aplicada, peritaje económico judicial y arbitral, y cuantificación de daños derivados de prácticas anticompetitivas y competencia desleal.

Ha participado como perito económico en procedimientos ante tribunales ordinarios, tribunales arbitrales y autoridades de competencia, abarcando casos de cárteles, abuso de posición dominante, competencia desleal, daños contractuales y disputas sobre valoración de activos. Su trabajo combina el rigor metodológico de la econometría moderna con la capacidad de comunicar resultados complejos de forma comprensible para jueces, árbitros y abogados.

Es autor de varios manuales de microeconometría aplicada, y compagina su actividad pericial con la docencia universitaria en econometría, análisis de datos y métodos cuantitativos para el Derecho.

Este manual refleja la confluencia de esas dos trayectorias: la académica, que exige precisión y transparencia metodológica, y la profesional, que demanda relevancia práctica y claridad expositiva ante tribunales.

\vfill

\begin{center}
\texttt{www.digitalreasons.es}
\end{center}

\clearpage
